% app-c-units — conventions, and the two traps that cost us real time
\section{Units and conventions}
\label{app:units}

Masses are quoted as $M/g$ unless stated otherwise; $m/g$ is the quark mass in
units of the coupling.  The natural combination in this paper's normalization
is $g^2 N/2\pi$, whereas 't~Hooft's classic solution uses $g^2 N/\pi$; the
large-$N$ comparison in Sec.~\ref{sec:modern} therefore plots
$(2\pi/N)^{1/2}\,M/g$ against $(2\pi/N)^{1/2}\,m/g$, exactly as Fig.~8(b) of
Ref.~\cite{Hornbostel:1988fb} did.  Spectra shown as functions of the coupling
use the compactified variable $\lambda = [1 + \pi m^2/g^2]^{-1/2}$, which maps
the free limit to $\lambda \to 0$ and the chiral limit to $\lambda = 1$, and
the dimensionless mass $M^2/(m^2 + g^2/\pi)$, both directly comparable to the
1990 figures.  In code units the resolution is $2K$ (integer, matching the
parity of the parton number) and single-parton momenta are odd integers
$k = 1, 3, \ldots$, so a structure function lives on $x = k/2K$ with measure
$dx = 1/K$.

Two conventions of Ref.~\cite{Hornbostel:1988fb} are traps for anyone
comparing against it, and we state them explicitly since both bit us.
First, \emph{the numbers printed in its Table~I are $M^2/(m^2+g^2/\pi)$,
although the column headers read $M/g$}.  The two agree at no coupling; the
identification was established by recomputing both quantities at the 1990
parameters and matching digit patterns, and our transcription of the table
carries the correction with full provenance.  Every comparison in this paper
converts to genuine $M/g$, and the cross-check that our conversion is right
--- recomputing the comparison in the table's native $M^2$ units --- is part
of the build (Appendix~\ref{app:data}).  Second, the parenthesized digits in
that table are \emph{not} error bars: they record the magnitude of the last
term retained in the extrapolation series, a convergence diagnostic the
authors themselves trusted only for $m/g \gtrsim 0.2$.  Where this paper
quotes an uncertainty, it is the full budget of Sec.~\ref{sec:budget}; where
it quotes the 1990 numbers, the parenthesis keeps its original meaning.

Captions in this paper state $N$, $B$, $2K$ (or the fit window), $m/g$, the
Fock-space truncation, and the Hamiltonian, for every panel; where the
original figure left a parameter unstated, the value we recovered and the
evidence for it are cited inline.
